% !TeX root = ../papers/nc-p3.tex
\section{二つの反対行列の階数二軌跡が有限な配置の全延長}
\label{sec:clean-configuration-extensions}
点スキームの同じ最初の配置を持つ窓を、他の成分からの延長も含めて比較する。
四辺有限な既知の窓を、共通の長さ二十のスキーム上での順序付き射影座標列
$(A,B,C,D)$ で記す。以下で同じ文字は標識付きの同じ射を表す。
\begin{theorem}\label{thm:clean-configuration-extensions}
標準正則窓 $(A,B,C,D)$ があり、$A$ 上で $C$ の平方を表す対称二次式行列と、
$B$ 上で $D$ の平方を表す対称二次式行列が、ともに階数2以下の軌跡を有限に持つとする。
最初の配置を $A$ に固定した全ての標準正則延長は次の四列で尽くされる。
\[
 (A,B,C,D),\qquad(A,D,C,B),\qquad
 (A,B,A,B),\qquad(A,D,A,D).
\]
初めの二列は常に正則である。後の二列はそれぞれ二つの隣接積の核を交互に用いて作り、
標準有限Frobenius検査と両対称判定（定理\ref{thm:bisymmetric-two-length-test}）を
通ったとき、かつそのときに限って採用する。この結論には点の被約性を仮定しない。
標識を忘れた添字保持同型は、$A$ の射影自己同型によってこのリストの順序付き列が移り合うことと同値である。
\end{theorem}
\begin{proof}
任意の標準正則延長は四辺有限となる。実際その第一射影が有限なので、
前節の系の証明で面を含む場合を除き、全射影の次元が1以下の場合の無分岐性と
グラフの同定を使えばよい。共通の有限スキームを $X$、その関数環を $H$ とする。
各文字束を自明化し、辺の完全交叉の完全跡を一つ固定する。
次の文字の平方空間は $(A^2)^\perp$ の単元倍であり、反対位置のものは $A^2$ の単元倍である
（定理\ref{thm:finite-square-chart}）。実際跡の取り直しは $H$ の単元倍で、
その単元は標数零の代数閉体上の有限環で平方根を持つ。
同じことが直前の文字にも逆辺で成り立つ。

$A^2$ の全階数の射影平方根は $A,C$ の二つに限られる。
定理\ref{thm:double-point-opposite-reconstruction}と
定理\ref{thm:nonreduced-opposite-root-separation}を $A$ 上の根 $C$ に使えばよい。
第二の定理の一般階数二の排除は、比較する別の平方根と $A$ が一つの窓の反対位置に
あることを必要としない。正規化後の十成分が全二次式の基底であり $X$ 上で階数1となるという条件で足りる。
$B^2$ 上では同じ推論で $B,D$ だけが残る。有限階数降下の仮定より $A,C$ は異なり、
$B,D$ も異なる。従って候補列は $(A,E,G,F)$ で
$E,F\in\{B,D\}$、$G\in\{A,C\}$ と書ける。

連続した三射が定まれば二辺の関係は隣接積の核として決まり、
その三次交叉の四次元基底が残る文字を除く全窓を決める
（定理\ref{thm:finite-square-chart}）。$G=C$ のとき $(A,E,C)$ を使うと、
$E=B$ なら元の窓、$E=D$ なら逆に読んだ窓に一致する。
逆転は二側の正則性・Koszul性を保存する。$E\ne F$ のときも
$(F,A,E)$ を使えば元の窓か逆転に一致し、$G=A$ とは両立しない。
従って残るのは $E=F$ と $G=A$ の両方を満たす二列である。
この二列の各辺は $R=\ker(A\otimes E\to H)$ とその転置である。
有限Frobenius窓の再構成後、両側の射影同定は両交換の対称性を与える
（定理\ref{thm:finite-opposite-product-gap}）。よって述べた判定で十分かつ必要である。
四移動は有限グラフ上で恒等なので、関係を保つ同型は全列に同じ $X$ の自己同型を誘導する。
逆にそのような自己同型で各射が射影同値になるなら全ての隣接積核が一致する。
\end{proof}

\begin{corollary}\label{cor:finite-reconstruction-one-edge}
上の仮定で一つの向き付き第一辺 $(A,E)$ を固定すると、異なる標準正則延長は高々二つである。
命題\ref{prop:finite-flip-obstructions}の $24\times24$ 行列の階数が20と異なればただ一つである。
一方の反対行列しか仮定できなくても、指定辺 $(A,B)$ と階数降下が有限な根 $C$ に対しては
全正則延長は高々二つである。非対称な第三座標は $I_{2X}$ の極小分解から決まり、
第三座標が $A$ の場合は反転辺の三次交叉から決まる。
特に十五次元成分の一般窓は、そのどの一辺を指定しても全成分の中で唯一の標準正則延長である。
最初の点配置だけを指定すれば、全標準正則延長は相異なる二つの向きで尽くされる。
\end{corollary}
\begin{proof}
一方向だけでも第三位置の候補は $A,C$ の二つであり、
どちらを選んでも固定した二辺の三次交叉から高々一つの窓を得る。
十五次元成分では系\ref{cor:normal-octics-general-finite-chart}により四位置とも仮定を満たす。
\path{computations/verify_clean_configuration_extensions.py} はその有理標本の四辺を両方向に読み、
八つの交代化行列の階数24を確認する。よって一般点でも反転辺による三次交叉は零である。
またその一般点の奇二配置は射影非同値なので、点配置のみを指定した二つの向きは同型にならない。
\end{proof}
